\section{Introduction}
\para{The access problem.} Breast cancer is the most common cancer in Indian women, with 192{,}020 new cases estimated for 2022 \cite{globocan2022}. Hospital series consistently report that a majority of patients are diagnosed at stage III or IV \cite{iips2023}. India's National Programme for Prevention and Control of Non-Communicable Diseases (NP-NCD) recommends clinical breast examination (CBE) every five years for women aged 30--65, delivered by health workers \cite{npncd}. A cluster-randomised trial of 151{,}538 women in Mumbai found that biennial CBE by trained primary health workers reduced breast-cancer mortality in women aged 50 and over \cite{mittra2021}. Uptake remains low: analyses of national surveys report that about 1\% of women aged 30--49 have ever had a CBE and 1.3\% of women aged 45 and over have had a mammogram \cite{bmcph2025}. Mammography needs fixed installations and radiologists. CBE detects palpable masses only.

\para{The role of handheld ultrasound.} Ultrasound images the breast without ionising radiation, works in dense breasts, and is now available as handheld probes connected to a phone or tablet. In a multicentre randomised trial in high-risk Chinese women, ultrasound detected all 14 cancers in the combined arm against 8 for mammography \cite{shen2015}. The bottleneck is interpretation: reading a whole-breast sweep requires a trained sonologist, and such readers are scarce in exactly the settings where CBE is the only screening available. Existing accessible tools in India address other modalities: iBreastExam measures tissue stiffness with a tactile sensor \cite{ibreastexam2022}, and Thermalytix analyses thermal images \cite{thermalytix2024}. Published AI work on handheld breast ultrasound mostly classifies lesions that a specialist has already found and framed \cite{handheldreview2025,berg2023}. Models trained on public breast-ultrasound datasets lose substantial accuracy when tested on a different dataset \cite{shiftbenchmark2026}.

\para{This work.} SweepScreen is a software pipeline that turns a health worker's handheld sweep into a per-woman screening decision and is designed for CPU-only deployment. Its contributions are:
\begin{enumerate}\itemsep1pt
\item An end-to-end architecture (Section~\ref{sec:arch}) that separates scan adequacy, frame-level lesion detection, lesion characterisation, temporal aggregation and threshold certification, with each stage independently testable and evaluated.
\item A unified data layer over five public datasets from Egypt, Brazil, Poland, Spain and China, with perceptual-hash deduplication and group-level leave-one-site-out splits (Section~\ref{sec:data}).
\item A frozen-foundation-model detector that trains on CPU, with augmentation that reproduces handheld-probe degradations, and a measured robustness profile under each degradation (Sections~\ref{sec:components} and~\ref{sec:results}).
\item A site-shift safety layer that sets the recall threshold at a new site from a small number of local cases, with exact finite-sample guarantees, and an evaluation of how often those guarantees hold on held-out sites (Section~\ref{sec:certify}).
\item A deployable service (HTTP API, health-worker web interface, versioned model registry, inference log, container) and a single script that regenerates every table, figure and number in this paper (Section~\ref{sec:deploy}).
\end{enumerate}
